% Section: limitations (\input from main.tex)
\section{Limitations and Threats to Validity}
\label{sec:threats}

\begin{itemize}
  \item \textbf{Platforms A and B share one chip, and B and C are VMs.} Platforms A and
        B both run on the same Apple M5, and Platforms B and C are virtual
        machines whose virtual CPUs are scheduled by a hypervisor we do not
        control. Bare-metal Platform~D (Intel Xeon W-2255, 10 cores / 20 threads)
        removes hypervisor noise, though it ran on a shared workstation with light
        background activity ($\text{load average} \approx 3.5$).
  \item \textbf{Different query sets.} \texttt{std::\allowbreak uniform\_int\_\allowbreak distribution}
        differs between libc++ and libstdc++, so Platform~A answered other
        random queries than B and C. Speed-ups are computed per platform and
        remain comparable, but absolute times and expansion counts of A should
        not be compared with B and C directly.
  \item \textbf{Short runs.} A Part~3 configuration (20 queries) takes
        19\,ms to 4.3\,s, so a single query often lasts only a few
        milliseconds. At this scale, OS noise and CPU
        frequency changes are visible. We report medians of 3 repetitions
        only, without spread.
  \item \textbf{Part~3 on Platforms~A and~C predates the final termination
        rule.} An earlier version of the shared-OPEN search stopped when OPEN
        was empty and a counter of busy threads was zero, read one after the
        other. That left a rare race, which we never observed in normal runs
        and could only trigger with injected delays; the \texttt{pending}
        counter of Section~\ref{sec:par-adapt} removes it. Platform~B was
        measured again with the final code: the times of the affected versions
        changed by about 3\,\% (median), as much as those of the versions whose
        code did not change, and every cost stayed optimal. The Part~3 times
        and test runs of Platforms~A and~C in this report were taken with the
        earlier version. The final code also passes \texttt{test\_par} on an
        x86-64 Linux VM of the same kind as Platform~C.
  \item \textbf{Thread start-up and short-run spread.} In the standard sweep, every
        parallel query creates and joins its $T$ threads and reports medians of 3
        repetitions without spread. The extended Platform~D suite (\texttt{run\_ext.sh})
        isolates search scaling via a persistent \texttt{ThreadPool}, pins threads to
        physical cores first, and reports 7-repetition standard deviation, IQR, and
        per-query $p_{95}$/$p_{99}$ latencies.
  \item \textbf{Graph built without OSMnx.} Because Nominatim was blocked
        and Overpass overloaded, the city graph comes from Geofabrik's extract
        with our own filter and simplification, cut by a bounding box rather
        than by the administrative boundary. It keeps service roads (except
        parking aisles and driveways), which OSMnx's \texttt{drive} network
        partly drops, so node counts differ from the lecturer's notebook.
  \item \textbf{Unequal sweeps.} Part~3 used up to $2\times$ the hardware
        threads, i.e.\ only $\{1,2,4\}$ on Platform~C, and left out the list
        versions with $h = 0$ because they would run for hours.
  \item \textbf{Synthetic parts of the workload.} The costly heuristic adds
        artificial work to model an expensive predictor, the dynamic and
        uncertain heuristics are noise models, and the congestion factors and
        traffic jams are invented, and the lecturer's speed profile is the same on
        every road. The road graph itself is real.
  \item \textbf{Memory is reclaimed late.} Removed list nodes are freed only
        when the list is destroyed. This is safe and fine for one search, but
        a long-running server would need hazard pointers or epoch-based
        reclamation.
  \item \textbf{No tuning sweep.} The backoff limits (1\,024 for locks, 256
        for validation retries) and the 3\% $f$-window were chosen once and not
        tuned per platform.
  \item \textbf{No thread pinning.} Threads were not pinned to cores, so on
        Platform~A we cannot tell how many of them ran on efficiency cores.
\end{itemize}
